%! AP Statistics — Unit 8, Lesson 8.6 — Homework
\documentclass[10pt]{article}
\usepackage{apstats-article}
\usepackage{apstats-boxes}
\newcommand{\TallMath}[1]{$\displaystyle #1\rule[-1.4em]{0pt}{3.2em}$}

\begin{document}

\pageheader{Unit 8, Lesson 8.6}{Homework}
\namedateperiod

\vspace{0.08in}

\textit{Show all work. Use all four SPDC steps for Problem 1.}

\begin{enumerate}[itemsep=10pt]

%% Problem 1 %%
\item \textbf{Sleep and Grade Level.}
A school counselor randomly selects 180 teenagers and classifies them by
\textbf{grade} (9th / 10th / 11th) and \textbf{sleep amount} ($<7$ hours / $7{+}$ hours per night).

\renewcommand{\arraystretch}{1.4}
\begin{center}
\begin{tabular}{@{} l c c c @{}}
\hline
           & \textbf{$<$7 hrs} & \textbf{7$+$ hrs} & \textbf{Total} \\
\hline
9th grade  & 30 & 30 & 60  \\
10th grade & 38 & 22 & 60  \\
11th grade & 42 & 18 & 60  \\
\hline
\textbf{Total} & 110 & 70 & 180 \\
\hline
\end{tabular}
\end{center}

A chi-square table excerpt (right-tail areas):

\renewcommand{\arraystretch}{1.3}
\begin{center}\small
\begin{tabular}{@{} c c c c c @{}}
\hline
\textbf{df} & $p = 0.10$ & $p = 0.05$ & $p = 0.025$ & $p = 0.01$ \\
\hline
1  & 2.706  & 3.841  & 5.024  & 6.635  \\
2  & 4.605  & 5.991  & 7.378  & 9.210  \\
3  & 6.251  & 7.815  & 9.348  & 11.345 \\
4  & 7.779  & 9.488  & 11.143 & 13.277 \\
\hline
\end{tabular}
\end{center}

\begin{enumerate}[label=\alph*., itemsep=6pt]
  \item \textbf{State.} Identify the test type. Write $H_0$ and $H_a$.

        Test type: \blank{5cm}

        $H_0$: \writeline
        $H_a$: \writeline

  \item \textbf{Plan.} Method: \blank{5.5cm}. Verify the Random condition: \writeline

        Compute expected counts (each row total = 60):

        $E_{\text{9th, }<7} = \dfrac{60 \times 110}{180} = $ \blank{2.5cm}
        \quad $E_{\text{9th, }7+} = \dfrac{60 \times 70}{180} = $ \blank{2.5cm}

        Are all expected counts $\ge 5$? \blank{2cm}

  \item \textbf{Do.} $\text{df} = ($ \blank{1.5cm} $-1)($ \blank{1.5cm} $-1) = $ \blank{2cm}

        Compute $\chi^2$ contributions and sum them below.

        \renewcommand{\arraystretch}{1.6}
        \begin{center}
        \begin{tabular}{@{} l c c @{}}
        \hline
                   & \textbf{$<$7 hrs} & \textbf{7$+$ hrs} \\
        \hline
        9th grade  & $\dfrac{(30 - \blank{2cm})^2}{\blank{2cm}} = $ \blank{2cm}
                   & $\dfrac{(30 - \blank{2cm})^2}{\blank{2cm}} = $ \blank{2cm} \\
        10th grade & \blank{4.5cm} & \blank{4.5cm} \\
        11th grade & \blank{4.5cm} & \blank{4.5cm} \\
        \hline
        \end{tabular}
        \end{center}

        $\chi^2 = $ \blank{4cm}

        Using the table above, the p-value is approximately \blank{3cm}.

  \item \textbf{Conclude.} At $\alpha = 0.05$:
        \writelines{2}
\end{enumerate}

\vspace{0.1in}

%% Problem 2 %%
\item \textbf{Degrees of Freedom.}
A student computes $\chi^2 = 9.21$ for a $2 \times 4$ contingency table.

\begin{enumerate}[label=\alph*.]
  \item What is the degrees of freedom? \blank{3cm}
  \item Using the table in Problem 1, what is the critical value at $\alpha = 0.05$?
        \blank{3cm}
  \item Based on $\chi^2 = 9.21$, can the student reject $H_0$ at $\alpha = 0.05$?
        Explain.
        \writeline
\end{enumerate}

\vspace{0.1in}

%% Problem 3 %%
\item \textbf{Conceptual Question.}
In what sense does a \emph{larger} $\chi^2$ value provide more evidence against $H_0$?
Use the ideas of expected counts, p-value, and the chi-square distribution in your answer.
\writelines{4}

\end{enumerate}

\end{document}
